\documentclass[10pt,a4paper]{amsart}
\usepackage[margin=19mm]{geometry}
\usepackage{amsmath,amssymb,mathtools}
\usepackage[scr=rsfs,cal=euler]{mathalfa}
\usepackage{xfrac}
\usepackage[unicode,hidelinks,bookmarks=false]{hyperref}
\newcommand{\ie}{i.e.\ }
\newcommand{\cf}{cf.\ }
\newcommand{\resp}{respectively\ }
\newcommand{\Subsection}{\subsection}
\providecommand{\nobreakdash}{\nobreak\hbox{-}}
\setlength{\emergencystretch}{2em}
\multlinegap0pt
\usepackage{amssymb}
\usepackage{xcolor}
\newtheorem{proposition}{Proposition}
\title{An upper bound for an exceptional automorphism group}
\author{Xu Zhuang}
\date{Source: arXiv:2609.03050v1 (CC BY 4.0)}
\begin{document}
\maketitle

\noindent\emph{Source and licence.} An upper bound for an exceptional automorphism group. Version arXiv:2609.03050v1 (CC BY 4.0), distributed by arXiv under the Creative Commons Attribution 4.0 International licence, http://creativecommons.org/licenses/by/4.0/, as recorded in the arXiv licence metadata for this exact version and shown on the abstract page https://arxiv.org/abs/2609.03050v1. Full licence notice: \texttt{licenses/arxiv-2609.03050-CC-BY-4.0.txt}.\par\smallskip

\noindent\textit{Editorial scope.} Counted unit: the article's proposition prop:unique (source0.tex lines 178--180). The article's complete original proof of it is delivered with it (source0.tex lines 182--190, one \texttt{proof} environment of 9 lines). No mathematics is written, repaired or re-derived: the statement and the proof are the article's own lines, in the article's own notation, and no retained line is substituted or re-flowed.  Retained uncounted as the source-local context the proof needs: lines 157--159; lines 164--166.  Nothing was omitted from any retained span, so the package carries no omission record.  No retained line contains a citation, so the wrapper carries no bibliography and no bibliography entry.  No retained line contains a figure environment, an image-inclusion command or drawing code, so no image is delivered and the package declares no asset dependency.  Editorial formatting only, in addition: an AMS \texttt{amsart} preamble that restates the article's own declarations and macros used by the retained lines, namely the environments \texttt{proposition} on separate counters, together with the standard packages those lines need.  The retained environments print in this document as Proposition 1 in order of appearance, whereas the article prints the same items with the numbering of its own full document.  The complete native source of the article as distributed in the arXiv e-print for this exact version is bundled unchanged as \texttt{sources/209258/source0.tex}.
\begin{equation}\label{eq:CS}
  g(L)\leq d_1g(F_1)+d_2g(F_2)+(d_1-1)(d_2-1).
\end{equation}

\begin{proposition}\label{prop:no-low-genus}
There is no degree-two morphism from $X$ to a curve of genus at most one.
\end{proposition}

\begin{proposition}\label{prop:unique}
If $v\in K\setminus k$ satisfies $[K:k(v)]=4$, then $k(v)=k(u)$. Equivalently, $k(u)$ is the unique rational subfield of $K$ having index four.
\end{proposition}

\begin{proof}
Let $B=k(u)$ and $B'=k(v)$, set $L=BB'$, and put $e=[K:L]$. Since $[K:B]=[K:B']=4$, the integer $e$ belongs to $\{1,2,4\}$.

If $e=1$, then $K=BB'$. Applying \eqref{eq:CS} to the two rational fields $B$ and $B'$ gives $g(X)\leq(4-1)^2=9$, contrary to $g(X)\geq11$.

Suppose that $e=2$. The tower law gives $[L:B]=[L:B']=2$, and $L=BB'$. Applying \eqref{eq:CS} inside $L$ gives $g(L)\leq(2-1)^2=1$. Since $[K:L]=2$, this contradicts Proposition~\ref{prop:no-low-genus}.

Therefore $e=4$. Since $B\subseteq L$ and both fields have index four in $K$, one has $L=B$. The same argument gives $L=B'$, and hence $B=B'$.
\end{proof}

\end{document}
